\chapter[(Аштавакра)]{Настанови для самореалізації (Аштавакра)}

[1.2] Щоб бути цілком вільним, уникай полону чуттєвих насолод так, як уникають отрути. Шукай терпимості, щирості, співчуття, простоти й чесності.

[1.3] Ти не земля, не вода, не вогонь і не повітря. Навіть не порожній простір. Звільнення --- це визнати себе чистою Свідомістю, свідком усього цього.

[1.4] Якщо ти зумієш побачити, що ти чиста Свідомість, відсторонившись від тіла, то вмить досягнеш щастя, спокою і звільнення.

[1.5] У тебе немає ні касти, ні обов'язків. Твоя природа невидима для очей, неприв'язана, безформна. Ти свідок усього. Будь щасливий.

[1.6] Правильне й хибне, насолода й біль існують лише в умі й тебе не стосуються. Не ти дієш і не ти страждаєш від наслідків, тож ти завжди вільний.

[1.7] Ти --- самотній Свідок і завжди вільний. Причина твого рабства --- вважати себе чимось понад це.

[1.8] Думка ``я дію'' --- це укус отруйної змії. Знання ``я нічого не роблю'' --- це мудрість-протиотрута. Будь щасливий.

[1.9] Саме лише розуміння ``я --- єдина Свідомість'' спопеляє все незнання й страждання у вогні одного мгновіння. Будь щасливий.

[1.10] Ти --- безмежна Свідомість, блаженство, найвищий екстаз, у якому всесвіт набуває форми, мов диво змії в мотузці. Будь щасливий.

[1.11] Те, що кажуть, --- правда: ``думка творить твою дійсність''. Якщо ти вважаєш себе обмеженим, то ти обмежений. Якщо вважаєш себе безмежним, то ти безмежний.

[1.12] Твоя справжня природа досконала, вільна, всюдисуща, єдина, відсторонена від усього, без бажань і в спокої. Сансара --- це ілюзія, породжена незнанням. Це існування, позбавлене самопізнання.

[1.13] Медитуй на себе як на саму лише Свідомість --- єдність як таку. Відкинь думку, що ти відокремлений, що свідомість твоя блукає. Відкинь думку, що є всередині й зовні.

[1.14] Довго ти перебував у пастці ототожнення з тілом, думаючи ``я --- людина''. Нехай знання ``я --- сама лише Свідомість'' стане мечем, що тебе звільняє, і будь щасливий.

[1.15] Споконвіку ти був вільним, світлим і незаплямованим. Заспокоювати ум силою --- це саме те, що тримає тебе в рабстві.

[1.16] Ти --- Свідомість, субстанція всесвіту. Дійсність існує завдяки тобі. Не тримай дрібного складу думок.

[1.17] Ти незумовлений, незмінний, безформний, нерухомий, незбагненний і спокійний. Не бажай нічого. Будь самою лише чистою свідомістю.

[1.18] Форма здається дійсною, тоді як безформне вічне. Усвідомивши це, ти вже не впадеш в ілюзію.

[1.19] Як у дзеркалі існують, зсередини й ззовні, відображення й відображений предмет, так найвище Я існує зсередини й ззовні тіла.

[1.20] Як усюдисущий простір існує зсередини й ззовні глечика, так позачасова єдність існує в цілокупності.
